\chapter{Experiments and Falsification}

\noindent\textbf{Status:} \texttt{REFERENCE-MODEL TESTS ACTIVE / PHYSICAL TESTS PREREGISTERED / EVIDENCE-CLASS SEPARATION REQUIRED}.

The experimental programme of Informational Dynamics of Time is organised as a sequence of increasingly demanding tests. The purpose of the sequence is to preserve a strict distinction between an algebraic identity, a deterministic reference-model reconstruction, a statistical anomaly, and a measured physical result. The temporal formalism therefore reaches experiment through explicit interfaces rather than through a single omnibus claim.

The canonical progression is
\[
\boxed{
\text{formal identity}
\longrightarrow
\text{reference implementation}
\longrightarrow
\text{withheld-variable test}
\longrightarrow
\text{preregistered physical test}
\longrightarrow
\text{independent replication}
}.
\]
Each promotion requires its own evidence receipt.

\section{Experimental observables of the temporal branch}

The temporal primitive supplies a positive activity measure
\[
d\Theta=\mathfrak a\,d\lambda,
\qquad \mathfrak a>0,
\]
and the local/reference comparison supplies the relational lapse
\[
N_R(x|r)=\frac{d\Theta_x}{d\Theta_r}
=\frac{\mathfrak a_x}{\mathfrak a_r}>0.
\]
Once a reference clock is physically calibrated by
\[
dt=T_r\,d\Theta_r,
\]
the local calibrated elapsed interval is
\[
d\widehat\tau_x=N_R(x|r)\,dt.
\]
This identifies the first experimental class: measurements of relative clock rate under independently characterised relational conditions. The primary observable is a dimensionless clock ratio, while conversion to seconds belongs to the calibration layer.

A second class concerns phase transport. On an admitted physical-clock calibration domain, the temporal branch is tested against the phase relation
\[
\Delta\phi_t=-\frac{E\Delta t}{\hbar}.
\]
The associated failure coordinate is direct: the temporal transport construction must recover the physical phase relation from its admitted bridge assumptions rather than by redefining the measured phase after the fact.

A third class concerns memory and retrodiction. Here the observable data are ordered checkpoints, event receipts, residence cells, winding increments, radial coordinates, and final retained states. These tests are especially useful because they permit exact withheld-lineage experiments before any physical retrocausal interpretation is considered.

\section{E003: withheld-lineage retrodiction}

The E003 protocol tests whether deliberately hidden lineage variables can be reconstructed from permitted checkpoints and declared model parameters. Its reference dynamics are
\[
X_{k+1}
=
\Phi_K(\Delta\tau_k;\mu_M)
\circ K_{u_k}(X_k),
\qquad
u_k=q_k\,\delta m_k,
\]
where the smooth Kepler-memory segment and the event kick are independently typed operations. The sealed truth variables are excluded from the estimator input.

The protocol contains several distinct identifiability regimes. If one factor of the event product is independently known, the complementary factor can be reconstructed from the inferred kick. If both \(q_k\) and \(\delta m_k\) are hidden, the event product remains the identifiable coordinate and factorisation is treated as a separate ambiguity class. Multi-event reconstruction is audited through the observation Jacobian
\[
J_R=\frac{\partial Y_C}{\partial z}.
\]
For a latent coordinate of dimension \(2N\), local observability requires
\[
\operatorname{rank}J_R=2N.
\]
A final-only multi-kick configuration supplies a mandatory underdetermined control.

The covariance-weighted residual is
\[
Q(z)=r(z)^T\Sigma_Y^{-1}r(z),
\]
and chronology controls are generated by checkpoint-block permutations with the same nominal information capacity. This separates reconstruction driven by ordered lineage from reconstruction achievable after chronology is destroyed.

The first reference run used three kicks, a six-dimensional latent coordinate, and a twelve-dimensional observation coordinate. The observability rank was six. The estimator converged in two iterations, with residual of order \(10^{-14}\) and maximum kick error of order \(10^{-14}\). Fifty randomised reference cases retained full rank and converged within the same numerical scale. The zero-kick and shuffled-chronology controls produced residuals many orders of magnitude larger than the exact reference solution.

The weighted diagnostic run again obtained rank six, with a modest condition number, while all tested nonidentity chronology permutations generated very large weighted residuals. The permutation ensemble in this diagnostic is an observability stress test rather than a sampling-based physical significance calculation.

The experimental meaning of E003 is therefore precise. It establishes that the admitted memory dynamics and observation map support exact reconstruction in the declared reference class, with controls for hidden-truth leakage, product ambiguity, rank deficiency, and chronology destruction. Its evidence class is \texttt{PASS\_MODEL\_INTERNAL} when the stated gates pass.

\section{Residence-bound compressed reconstruction}

The retrodiction carrier has been reduced from a full Cartesian packet to a winding--radius packet on each admitted active-attractor stratum. With
\[
\rho_k=\|r_k-c_{a_k}\|>0,
\]
and ordered winding increments \(\Delta W_k\), the pre-final positions are reconstructed by
\[
r_k=c_{a_k}+\rho_k
\begin{pmatrix}
\cos(\theta_{k-1}^{(a_k)}+2\pi\Delta W_k)\\
\sin(\theta_{k-1}^{(a_k)}+2\pi\Delta W_k)
\end{pmatrix}.
\]
The radial packet contains \(N-1\) new scalar coordinates in place of the \(2N-2\) Cartesian coordinates of the baseline carrier. Thus
\[
\boxed{
\frac{N_{\rm radial}}{N_{\rm Cartesian}}=\frac12
}
\]
for \(N>1\).

The 07V residence binding attaches every retained radius to the checkpoint index, active attractor, exact binary64 radius, source residence cell, post-segment state commitment, and preceding radial commitment. The content-addressed chain therefore supplies provenance for the compressed coordinate used by the decoder. The hosted reference suite for this gate passed 557 tests. The broader global injectivity coordinate remains a separate frontier beyond the declared positive-radius stratified domain.

\section{E004: prospective retrocausal discrimination}

Chapter~9 defines the preregistration boundary for a genuinely retrocausal candidate. The primary statistical target is
\[
P(X_{\mathrm{past}}\mid Y_{\mathrm{future}})
\stackrel{?}{\neq}
P(X_{\mathrm{past}}),
\]
with the future condition selected after the earlier measurement event and with the complete analysis contract committed beforehand.

Admission of a physical candidate requires, at minimum, cryptographic timestamp commitments, blinded future-condition assignment, a classical-channel audit, shared-clock and scheduler controls, sham/null trials, a fixed exclusion policy, and replication. The raw earlier record is immutable. Analysis proceeds from committed data products and an analysis plan whose hash predates condition unblinding.

The result ladder used throughout this programme is
\[
\texttt{RAW\_OBSERVATION}
\rightarrow
\texttt{STATISTICAL\_EFFECT}
\rightarrow
\texttt{CLASSICAL\_CHANNEL\_AUDIT}
\rightarrow
\texttt{PHYSICAL\_CLAIM\_STATUS}.
\]
A statistical deviation is therefore only one coordinate of the experimental conclusion.

\section{Falsification matrix}

The experimental programme contains five principal falsification families.

\begin{enumerate}
\item \textbf{Temporal phase transport.} Test recovery of \(\Delta\phi_t=-E\Delta t/\hbar\) on the declared calibration domain.
\item \textbf{Spacetime closure.} Combine the independently developed temporal and spatial branches and test the reconstructed phase against
\[
\Delta\phi=\frac{1}{\hbar}
\left(\mathbf p\cdot\Delta\mathbf x-E\Delta t\right).
\]
\item \textbf{NOW and bifurcation.} Test branch integrity, realised-lineage uniqueness, reconstructive reversibility where specified, and repeated-state re-entry consistency.
\item \textbf{Retrodiction and retrocausal discrimination.} Use E003 for model-internal reconstruction and E004 for the preregistered physical candidate test.
\item \textbf{Bell/CHSH discrimination where the experimental topology requires it.} For independently selected settings \(a,a',b,b'\), evaluate
\[
S=E(a,b)+E(a,b')+E(a',b)-E(a',b'),
\]
with explicit timing, selection, detection, and communication controls.
\end{enumerate}

\section{Evidence classes and provenance}

The canonical verdict classes are
\[
\begin{gathered}
\texttt{PASS\_MODEL\_INTERNAL},\quad
\texttt{FAIL\_MODEL\_INTERNAL},\quad
\texttt{CLASSICAL\_COMPATIBLE},\\
\texttt{ANOMALY\_CANDIDATE},\quad
\texttt{RETROCAUSAL\_CANDIDATE},\quad
\texttt{NONCLASSICAL\_CORRELATION\_CANDIDATE},\\
\texttt{MEASURED\_PHYSICAL\_RESULT}.
\end{gathered}
\]
Every promoted result retains the raw-data identity, source commit, parameter manifest, test coordinate, uncertainty model, and failure receipt. This provenance rule is part of the scientific content of the programme: a temporal claim is evaluated together with the chronology by which its evidence was produced.

The resulting experimental architecture mirrors the theory itself. Temporal order generates an observable lineage; memory makes that lineage addressable; retrodiction tests whether the retained structure is sufficient to reconstruct hidden causes; preregistration then supplies the boundary at which a future-conditioned physical effect can be distinguished from reconstruction, leakage, scheduling, and ordinary forward causation.
